\begin{Solbook}{1}
 \par \noindent  a) Vô nghiệm \hspace {2cm} b) $x=10, y=7$.  \hfill \qedEX 
\end{Solbook}
\begin{Solbook}{2}
 \par \noindent  a) $x>2$ \hspace {2cm} b) $4$ \hfill \qedEX 
\end{Solbook}
\begin{Solbook}{3}
 \par \noindent Áo khoác: $300$ chiếc; Quần jean: $200$ chiếc. \hfill \qedEX 
\end{Solbook}
\begin{Solbook}{4}
 \par \noindent Điểm $C$ cách điểm $A$ khoảng $179{,}3$ m. \hfill \qedEX 
\end{Solbook}
\begin{Solbook}{5}
 \par \noindent \immini {\begin {enumerate} \item Chứng minh các tam giác $OMA$ và $OPB$ bằng nhau để suy ra $OM=OP$. Tam giác $NMP$ có $NO$ vừa là phân giác vừa là đường cao nên chứng minh được tam giác $NMP$ cân. \item Chứng minh $\triangle OMN=\triangle ONP$, từ đó chứng minh $\triangle OBP=OIP$. Từ đó suy ra $OI=OB=R$ và $MN$ là tiếp tuyến của $(O)$. \item Chứng minh $\triangle OAM\backsim \triangle NBO$, từ đó suy ra đẳng thức tỉ lệ và suy ra điều phải chứng minh. \end {enumerate}} {\begin {tikzpicture}[scale=0.7, font=\footnotesize ] \path (0:0) coordinate (O) --(0:2.5) coordinate (B) --([turn]-90:3) coordinate (Bt) ; \path (0:0) coordinate (O) --(180:2.5) coordinate (A) --([turn]90:3) coordinate (At) ; \coordinate (M) at (-2.5,2); \coordinate (P) at (2.5,-2); \path [name path=at] (At)--($(At)!2!(A)$); \path [name path=bt] (Bt)--($(Bt)!3!(B)$); \coordinate (M') at ($(O)!2!-90:(M)$); \path [name path=om] (O)--(M'); \path [name intersections={of=om and bt,by={N}}]; \coordinate (I) at ($(M)!(O)!(N)$); \draw (O) circle (2.5) (M)--(P) (O)--(N)--(M)--cycle (O)--(I) ; \draw (A)--(B) (At)--($(At)!2.25!(A)$) node[left]{$d$} (Bt)--($(Bt)!2.25!(B)$) node[right]{$d'$}; \foreach \x /\g in{A/-135,B/-45,O/-90,M/-180,P/0,N/0,I/90} \fill (\x ) circle (1.5pt) +(\g :.3) node{$\x $}; \end {tikzpicture}}  \hfill \qedEX 
\end{Solbook}
\begin{Solbook}{6}
 \par \noindent  Đặt $R=2024$. Gọi $H,K$ lần lượt là giao điểm của các đường vuông góc từ $O_1,O_2$ xuống $d$ với đường thẳng qua $(I)$ song song với $d$ ($I$ là tâm của đường tròn cần tìm bán kính). Khi đó nếu gọi $r$ là bán kính của đường tròn nhỏ thì tam giác $IO_1H$ vuông tại $H$ có $IO_1=R+r,OH=R-r,HI=R$. Từ đó thiết lập phương trình theo Định lý Pytagore và tính được $r=\dfrac {1}{4}R=506$.  \hfill \qedEX 
\end{Solbook}
